\section{Estimates for class F}\label{app:forward-estimates}

This appendix proves the scalar inequalities of \cref{sec:forward}, in its notation: $\nu=4/a^2$,
the ratios $z,w,\hat p,\hat d\ge1$, $\hat s_\pm$, $\hat K$, $\hat K_\pm$, the factors $h,q,k$, the
row coefficients $A,B,\eta,n_1,n_2$ and the floor $F_x(r)$ of \cref{prop:fwd-window}. Throughout,
\[
f(r)=\frac{r-r^{-1}}2-\log r,\qquad \beta(y)=\frac{y^2+y+1}{(1+y)^2}=1-\frac1{1+y}+\frac1{(1+y)^2},
\]
so that $\beta(w)$ is the value $1-x+x^2$ of the floor $F_x$ at $r=1$ for $x=w/(1+w)$.

\subsection{The edge function}\label{app:fwd-edge}

Let $G$ be the edge function of \eqref{eq:fwd-cost}. Its partial derivatives are
\begin{align*}
\partial_bG(a,b)&=3a-2b\log(a+b)-\frac{a^2+b^2}{a+b}+b\log(b^2+4)+\frac{b^3}{b^2+4},\\
\partial_a\partial_bG(a,b)&=\frac{2(a^2+ab+b^2)}{(a+b)^2},\qquad
\partial_a\partial_b^2G(a,b)=\frac{2a(b-a)}{(a+b)^3}.
\end{align*}
\begin{enumerate}
\item (Monotonicity) The bound $\log y\ge1-1/y$ at $y=(b^2+4)/(a+b)^2$ gives
  $\partial_bG(a,b)\ge a(a+2b)(b^2-ab+8)/\bigl((a+b)(b^2+4)\bigr)$, which is positive when
  $ab\le8$. So $G$ increases in each argument while the product of its arguments is at most $8$.
  All products of consecutive active margins are below $4$.
\item (Rectangle inequality) Since $\partial_bG(x,\cdot)-\partial_bG(a,\cdot)\ge0$ for $x\ge a$,
  integrating in the second variable gives $G(x,y)\ge G(x,a)+G(a,y)-G(a,a)$ for $x,y\ge a>0$.
\item (Concavity) If $a^2\ge4/3$, then $G(a,\cdot)$ is concave on $[a,4/a]$. Since
  $\partial_a\partial_b^2G\ge0$ for $a\le b$, it suffices to check $\partial_b^2G<0$ at $a=b$ when
  $4/3\le b^2\le4$, and at $a=4/b$. With $\varsigma=4b^2/(b^2+4)\in[1,2]$ one has
  $\partial_b^2G(b,b)=-\log\varsigma-\frac32+\frac54\varsigma-\frac18\varsigma^2$, which increases in
  $\varsigma$ (its derivative is $(\varsigma-1)(4-\varsigma)/(4\varsigma)$) up to the value
  $\frac12-\log2<0$ at $\varsigma=2$; and $\partial_b^2G(4/b,b)=\log\varrho+1-\varrho<0$ with
  $\varrho=b^2/(b^2+4)$.
\item (Scaling) $G(ax,ay)=a^2g_\nu(x,y)$ with
  $g_\nu(x,y)=3xy-\nu-(x^2+y^2)\log(x+y)+\frac12x^2\log(x^2+\nu)+\frac12y^2\log(y^2+\nu)$; we write
  $g_\nu(x)=g_\nu(1,x)$. In particular $g_\nu(1)=3-\nu-\log\frac4{1+\nu}\ge\nu(3-\nu)/(1+\nu)\ge0$
  for $0<\nu\le3$, by $\log y\ge1-1/y$ again; so $G(a,a)\ge0$ when $a^2\ge4/3$.
\end{enumerate}

\subsection{Logarithm bounds}\label{app:fwd-logs}

All logarithms are bounded by rational functions through the following elementary inequalities,
each proved from the sign of a derivative. For $v\ge1$ and $\tau=(v-1)/(v+1)$,
\[
2\tau+\tfrac23\tau^3\le\log v\le2\tau+\frac{2\tau^3}{3(1-\tau^2)},
\]
and for $h\ge1$, $0\le y<1$,
\[
\frac{(h-1)^3}{3h(h+1)}\le f(h)\le\frac{(h-1)^3}{6},\qquad
\log(1-y)\le-y-\frac{y^2}2-\frac{y^3}3 .
\]
The lower bound for $f$ holds because the difference vanishes at $h=1$ and has derivative
$(h-1)^4/(6h^2(h+1)^2)$; it is the Pad\'e bound, used for every $\log h$ below and, in the form
$\log v\le\frac{v-v^{-1}}2-\frac{(v-1)^3}{3v(v+1)}$, for other logarithms. The upper bound for
$\log v$ at $v=5/4$ gives $\log\frac54<\frac9{40}$.

\subsection{The threshold and the budget}\label{app:fwd-budget}

\runin{The threshold.} With $a_0^2=20/11$ and $\nu_0=11/5$,
$g_{\nu_0}(1)=3-\frac{11}5-\log\frac{20}{16}=\frac45-\log\frac54$, so
$G(a_0,a_0)=\frac{20}{11}\bigl(\frac45-\log\frac54\bigr)>\frac{20}{11}\cdot\frac{23}{40}=\frac{23}{22}$.
With $\sqrt3>1.732$ and $\Aoct<3.1268$ this gives $\frac{7\sqrt3}4G(a_0,a_0)>\Aoct$, which proves
\cref{lem:fwd-threshold}.

\runin{The budget.} By the rectangle inequality the seven-entry bound becomes
$\frac{\sqrt3}{\nu}\bigl(2g_\nu(\hat p)+2g_\nu(z)+2g_\nu(w)+2g_\nu(\hat d)-g_\nu(1)\bigr)\le\Aoct$
(\cref{sec:forward}, Step~3). The four ratios lie in $[1,\nu]$, where $x\mapsto g_\nu(x)$ is
increasing and concave by \cref{app:fwd-edge}. For an increasing concave function $\phi$ on $[1,M]$
and four points of $[1,M]$ with sum at least $\sigma$, where $1<\sigma-3\le M$,
\[
\phi(x_1)+\phi(x_2)+\phi(x_3)+\phi(x_4)\ge3\phi(1)+\phi(\sigma-3):
\]
lower the points to sum exactly $\sigma$, and use that a concave function lies above the chord from
$1$ to $\sigma-3$. If $\hat p+z+w+\hat d\ge(5\nu+1)/3$, this with $x_*=(5\nu-8)/3\in(1,\nu]$ bounds
the bracket below by $5g_\nu(1)+2g_\nu(x_*)$. Write
\[
g_\nu(x)=3x-\nu-\tfrac12\log\frac{(1+x)^2}{1+\nu}-\tfrac{x^2}2\log\frac{(1+x)^2}{x^2+\nu},
\]
and bound both logarithms above by
$L_+(N,M)=2\frac{N-M}{N+M}+\frac{(N-M)^3}{6NM(N+M)}$, the upper bound of \cref{app:fwd-logs} at
$v=N/M\ge1$. This gives a rational lower bound $\bar g_\nu(x)$, and the remaining claim is
\[
5\bar g_\nu(1)+2\bar g_\nu(x_*)>\frac{2501}{1384}\,\nu\qquad(11/5\le\nu\le3).
\]
After multiplication by the positive denominator
$934200(\nu-1)^2(\nu+1)(\nu+5)(25\nu^2-71\nu+64)(25\nu^2-41\nu+34)(50\nu^2-121\nu+89)$, this is the
positivity on $[11/5,3]$ of a polynomial of degree~$11$ in $\nu$, which we verify by expanding it
in the Bernstein basis of $[11/5,3]$: all twelve coefficients are positive (this is the argument of
\path{GLOBAL_FORWARD_LINEAR_BUDGET} in~\cite{Notes}). Since $\sqrt3>173/100$ and
$\frac{173}{100}\cdot\frac{2501}{1384}=\frac{2501}{800}>\Aoct$, the cost would exceed $\Aoct$. This
proves \cref{lem:fwd-budget}.

\subsection{The centered system}\label{app:fwd-centered}

The following identities hold for $\nu,z,w,h>0$, $h\ne1$; each is an equality of rational
expressions in the variables and the logarithms.
\begin{enumerate}
\item (Floors) $F_{z/\hat s_-}(h)=\frac34+L(z,h)$ with
  \[
  L(z,h)=\frac{(z-1)\bigl(z-1+2z(h-1)\bigr)}{4\hat s_-^2}+\frac{zh}{\hat s_-}\cdot\frac{f(h)}{h-1},
  \]
  and $L(z,h)\ge0$ for $z\ge1$, $h>1$.
\item (Central row) $\frac34(\hat s_-+\hat s_+)-A=e_0-\varphi$.
\item (Left row) With $\zeta=z-1$ and $H=h-1$, the quantity $r=\eta A-B-\frac34(n_1+n_2)$ of
  \eqref{eq:fwd-centered} is given by
  \begin{align*}
  8\hat s_-\hat K_-\,r={}&(3-\nu)\zeta(3-\nu+2\zeta)\\
  &+H\bigl((3-\nu)^2(\zeta^2+4\zeta+2)+8(3-\nu)\zeta z\bigr)
  +2H\zeta^2(8+4\zeta-\zeta^2)\\
  &+\frac{4\hat s_-H^2z(3-\nu+4\zeta+\zeta^2)}{h}\\
  &+8\hat s_-\hat K(1+\hat s_-H)\,\varphi-8\hat s_-\hat K_-\,zf(h).
  \end{align*}
  It depends on $w$ only through $\varphi$.
\item (Coefficients) $\hat s_-\hat K_-(n_2-n_1)=\hat K(\hat s_-\hat s_+-\hat K)
  +H\bigl(\hat s_-^4-\hat K^2+\hat K\hat s_-^2(\hat s_+-\hat s_-)\bigr)$, which is positive for
  $\nu<3$, $1\le z\le w$, $h\ge1$, because $\hat s_-\hat s_+\ge4>\hat K$ and
  $\hat s_-^4\ge16>\hat K^2$.
\end{enumerate}
Two monotonicity facts complete the setup. The floor increases in the factor: for $x\ge1/2$,
$\partial_rF_x(r)=x^2-x\,(r-1-\log r)/(r-1)^2\ge x^2-x/2\ge0$, since
$r-1-\log r\le(r-1)^2/2$. And in the variable $w$, with $\nu,z,h$ fixed,
\[
n_2(w)-n_2(z)=\eta\,(w-z),\qquad A(w)-A(z)\le(w-z)\,\beta(w)\qquad(z\le w),
\]
because $\partial_wA=\beta(w)$ and $\beta'(y)=(y-1)/(1+y)^3\ge0$ for $y\ge1$.

\subsection{The three residual signs}\label{app:fwd-signs}

\runin{Lowering the larger neighbour.} The residuals of \cref{sec:forward} depend on $w$ through
$A$, $n_2$, $\varphi$ and the opposite floor. For $\mathcal E_v=r-n_1L(z,h)-n_2/64$ the increments
above give
$\mathcal E_v(w)-\mathcal E_v(z)\le\eta(w-z)\bigl(\beta(w)-\frac{49}{64}\bigr)\le0$, since
$\beta(w)\le\beta(4/3)=37/49<49/64$. For
$\mathcal E_\Sigma=r-n_1/35-(n_2-n_1)L(w,\hat s_+^2/\hat K_+)$ put $V(w)=\frac34+L(w,\hat s_+^2/\hat K_+)$,
the floor at the least admissible factor; then $V(w)\ge\beta(w)$, and $V$ is nondecreasing in $w$
(the coefficients of $L$ increase in $w$, $hf(h)/(h-1)$ increases in $h$, and
$\hat s_+^2/\hat K_+$ increases in $w$), so
\[
\mathcal E_\Sigma(w)-\mathcal E_\Sigma(z)
\le\eta(w-z)\bigl(\beta(w)-V(w)\bigr)-(n_2(z)-n_1)\bigl(V(w)-V(z)\bigr)\le0 .
\]
The neighbour residual is lowered in the same way to $w=\max(z,4/3)$. After lowering, $w=z$ for
$\mathcal E_\Sigma$ and $\mathcal E_v$, and the factor ranges over the symmetric domain
\begin{equation}\label{eq:fwd-symdomain}
1\le z\le\tfrac43,\quad\ 1+\tfrac65z\le\nu<3,\quad\
\frac{(1+z)^2}{z^2+\nu}\le h\le\frac{(1+z)(5\nu-2-3z)}{3(z^2+\nu)};
\end{equation}
the lower bound for $\nu$ is exactly the condition that the interval for $h$ be nonempty.

\runin{Concavity in the factor.} As functions of $h$, with the other variables fixed, the three
residuals have the form $c_0+c_1h+c_2h^2+c_{-1}/h+(d_0+d_1h)\log h$, with second derivative
$(2c_2h^3+2c_{-1}+d_1h^2-d_0h)/h^3$. For $\mathcal E_v$ and $\mathcal E_N$ it equals
$-zP(h)/(\hat K_-h^3)$ with
\[
P(h)=(\nu-z)h^2(2zh-\hat s_-)+\hat K_-h-2z\hat s_- .
\]
For $h\ge1$, $P'(h)=2(\nu-z)h(3zh-\hat s_-)+\hat K_->0$, and at the least factor
$H_0=\hat s_-^2/\hat K_-$,
\begin{gather*}
P(H_0)=-\frac{\hat s_-(2z+1-\nu)}{\hat K_-^3}\,Q(\nu,z),\\
Q(\nu,z)=(1-z)\nu^2-z(z^3+7z^2+5z+3)\nu+z^2(z^4+2z^3+6z^2+5z+2).
\end{gather*}
On the domains of these residuals, $1\le z\le5/3$, $\nu<3\le2z+1$ and $\nu\ge1+6z/5$. The
polynomial $Q$ decreases in $\nu$, and at $\nu=1+6z/5$, with $x=z-1\in[0,\frac23]$,
\[
Q=x^6+\tfrac{34}5x^5+\tfrac{78}5x^4+\tfrac{124}{25}x^3-\tfrac{832}{25}x^2-\tfrac{1216}{25}x-\tfrac{96}5
\le-\tfrac{42176}{2025}x^2-\tfrac{1216}{25}x-\tfrac{96}5<0,
\]
bounding the four positive terms by
$\bigl((\frac23)^4+\frac{34}5(\frac23)^3+\frac{78}5(\frac23)^2+\frac{124}{25}\cdot\frac23\bigr)x^2$.
So $P\ge0$ on $[H_0,\infty)$, and $\mathcal E_v,\mathcal E_N$ are concave in $h$ there. A concave
function on an interval lies below its tangent at an endpoint, and the tangent is affine, so the
residual is negative on the interval as soon as the tangent is negative at both endpoints. This
leaves two values for each residual: the value at one endpoint and the tangent's value at the
other. The sum residual is treated in the variable $\tau=(h-1)/h$ instead (below).

\runin{The endpoint values.} On \eqref{eq:fwd-symdomain} we use the coordinates
\[
\alpha=\frac{z-1}{z+1}\in[0,\tfrac17],\qquad t=\frac{(1+z)^2}{z^2+\nu}-1\in(0,\tfrac14],
\]
in which $\hat K/\hat K_-=1-\alpha(1+t)$ and $\hat s_-^2/\hat K_-=1+t$; the upper bound for $t$
follows from $5(z^2+1+\frac65z)-4(z+1)^2=(z-1)^2$. The endpoint values involve, besides rational
functions of $(\alpha,t)$, the two logarithms $\log(1+t)$ and $\log(1-\alpha(1+t))$, each with a
coefficient that is a polynomial in $(\alpha,t)$ divided by a positive one.
\begin{itemize}
\item $\mathcal E_v<0$: the tangent is taken at $h=1+t$, so the two values are
  $\mathcal E_v(1+t)$ and the tangent's value at the cost endpoint.
\item $\mathcal E_\Sigma<0$: here $\mathcal E_\Sigma=A_0+A_1H+z\bigl(DH^2/(2h)-f(h)\bigr)$ with
  $H=h-1$, $D=\alpha+t+\alpha t$ and coefficients $A_0,A_1$ free of $h$. The Pad\'e bound for
  $f(h)$ gives $\mathcal E_\Sigma/H\le\Psi(\tau)$ with $\tau=H/h$, where
  $\Psi(\tau)=A_0(1/\tau-1)+A_1+z\bigl(D\tau/2-\tau^2/(3(2-\tau))\bigr)$ is concave once
  $A_0\le0$. The tangent at the largest admissible value
  $\tau_c=\frac{2-2\alpha^2-t(5+3\alpha+2\alpha^2)}{5-3\alpha-2\alpha^2-t(5+3\alpha+2\alpha^2)}$
  gives the two values $\Psi(\tau_c)$ and $\Psi(\tau_c)+\Psi'(\tau_c)(\tau_I-\tau_c)$, where
  $\tau_I=t/(1+t)$ is the least value. The domain also satisfies $t\ge39\alpha/40$.
\item $\mathcal E_N<0$ for $w\ge4/3$: after lowering $w$ to $\max(z,4/3)$ there are two faces. On
  the symmetric face $w=z\in[\frac43,\frac53]$ one has $\alpha\in[\frac17,\frac14]$ and
  $t\in[\frac{13}{100},\frac14]$, with the same two logarithms. On the asymmetric face $w=4/3\ge z$
  the budget gives $\nu\ge\frac95+\frac35z\ge\frac{12}5$, the variables are
  $(\nu,z)\in[\frac{12}5,3]\times[1,\frac43]$, and the logarithms are $\log(\hat s_-^2/\hat K_-)$,
  $\log q$ and $\log(\hat s_+^2/\hat K_+)$. On each face the tangent is taken at
  $h=\hat s_-^2/\hat K_-$.
\end{itemize}
Each logarithm is replaced by its Pad\'e upper bound, by the cubic bound for $\log(1-y)$, or by the
two-term lower bound of \cref{app:fwd-logs}, according to the sign of its coefficient. Each value
then becomes a rational function with positive denominator, and the claim is that its numerator,
a polynomial in two variables of degree at most~$14$, is negative. The polynomial signs used are:
for $\mathcal E_v$, the two logarithm coefficients and the two numerators on
$[0,\frac17]\times[0,\frac14]$; for $\mathcal E_\Sigma$, the coefficients and numerators of $A_0$
(which is nonpositive) and of the first value on $[0,\frac17]\times[0,\frac14]$, and those of the
second value on the boxes $[0,\frac1{10}]\times[0,\frac14]$ and
$[\frac1{10},\frac17]\times[\frac{39}{400},\frac14]$, which cover the part of the domain with
$t\ge39\alpha/40$; for $\mathcal E_N$, the coefficients and numerators of both values on
$[\frac17,\frac14]\times[\frac{13}{100},\frac14]$ (symmetric face) and on
$[\frac{12}5,3]\times[1,\frac43]$ (asymmetric face). We verify each of these polynomial signs by
expanding the polynomial in the Bernstein basis of its box: all coefficients have the required
sign. Where a numerator vanishes at a corner of its box and the inequality must be strict, the
expansion is applied to the numerator plus a small linear form that is positive away from that
corner. Another proof of the three signs, by displayed bounds on the coefficients of the numerators
in one variable at a time, an integral bound for the neighbour residual and a split of the sum
residual at $D=38/243$, is given in \path{HAND_FORWARD_LARGER_NORM_HAND_PROOF},
\path{HAND_FORWARD_SYMMETRIC_SUM_REDUCTION}, \path{HAND_FORWARD_SYMMETRIC_SUM_HIGH_D},
\path{HAND_FORWARD_NEIGHBOR_INTEGRAL} and \path{HAND_FORWARD_ASYMMETRIC_NEIGHBOR_BOUND}
of~\cite{Notes}.

\subsection{The conditional Casimir bound}\label{app:fwd-casimir}
Assume the hypotheses of \cref{lem:fwd-casimir}. Then $0\le\varphi\le1/100$ and
$\frac34\le e_0\le\frac{73}{80}$, so $0\le E\le e_0+\hat s_-\delta u+\hat s_+\delta v$. The Casimir is
$\hat K^2\bigl(4uv-E^2\bigr)/(16\nu^2)$ and decreases in $E\ge0$, so it is at least
$\hat K^2Q(\delta u,\delta v)/(16\nu^2)$ with
\[
Q(x,y)=4\Bigl(\frac34+x\Bigr)\Bigl(\frac34+y\Bigr)-\bigl(e_0+\hat s_-x+\hat s_+y\bigr)^2 .
\]
Write $c=1/35$ and $b=1/64$. The function $Q$ does not increase in $x$ along the norm polygon: for
$x\le c-y$,
\[
Q(c-y,y)-Q(x,y)=(c-y-x)\bigl(3+4y-\hat s_-(2e_0+\hat s_-(x+c-y)+2\hat s_+y)\bigr)\le0,
\]
since $2\hat s_-e_0\ge3$ and $\hat s_-\hat s_+\ge4$. Along the edge $x+y=c$ it is concave, by the
exact identity
\[
Q(c-y,y)-\bigl((1-64y)\,Q(c,0)+64y\,Q(c-b,b)\bigr)=\bigl(4+(\hat s_+-\hat s_-)^2\bigr)\,y\,(b-y),
\]
which is nonnegative for $0\le y\le b$. So it suffices to bound $\hat K^2Q/(16\nu^2)$ at the two
vertices $A=(c,0)$ and $B=(c-b,b)$. The ordering and the budget give
$\hat s_-\le\min\bigl(\frac73,\frac{5\nu+1}6\bigr)$ and $\hat s_+\le\frac73$, and
$\hat Ke_0=\hat K-\hat K^2/8+1$. With these, the two vertex values are bounded below by explicit
functions of $\nu$ alone:
\[
\phi_A(\nu)=\frac{\hat K^2}{16\nu^2}\Bigl(\frac{327}{140}-\bigl(e_0+\tfrac1{35}\min(\tfrac73,\tfrac{5\nu+1}6)\bigr)^2\Bigr),
\]
and $\phi_B$, defined in the same way at $B$ with $\hat s_-(c-b)+\hat s_+b$ bounded by
$(c-b)(\hat s_-+\hat s_+)+(2b-c)\hat s_+\le(c-b)\bigl(2+\frac53(\nu-1)\bigr)+\frac73(2b-c)$ for
$\nu\le13/5$ and by $\frac73c$ for $\nu\ge13/5$. Multiplied by $16\nu^2$, each of the inequalities
$\phi_A>23/125$ and $\phi_B>23/125$, on $[\frac{11}5,\frac{13}5]$ and on $[\frac{13}5,3]$, is the
positivity of a quartic polynomial. In the variable $y\in[0,1]$ of the interval
($\nu=\frac{11}5+\frac25y$, resp.\ $\nu=\frac{13}5+\frac25y$) the four quartics are
\begin{align*}
&\tfrac{6283}{153125}-\tfrac{92}{3125}y-\tfrac{20162}{1378125}y^2+\tfrac{34}{7875}y^3
  -\tfrac{289}{1102500}y^4 &&(\phi_A,\ [\tfrac{11}5,\tfrac{13}5]),\\
&\tfrac{89}{87500}+\tfrac{11419}{65625}y+\tfrac{10643}{393750}y^2+\tfrac1{375}y^3
  -\tfrac1{2500}y^4 &&(\phi_A,\ [\tfrac{13}5,3]),\\
&\tfrac{9671}{306250}-\tfrac{827}{65625}y-\tfrac{2033761}{176400000}y^2
  +\tfrac{21823}{5040000}y^3-\tfrac{19321}{70560000}y^4 &&(\phi_B,\ [\tfrac{11}5,\tfrac{13}5]),\\
&\tfrac{9203}{800000}+\tfrac{211601}{1200000}y+\tfrac{195547}{7200000}y^2+\tfrac1{375}y^3
  -\tfrac1{2500}y^4 &&(\phi_B,\ [\tfrac{13}5,3]).
\end{align*}
On $[\frac{13}5,3]$ the only negative term $-y^4/2500$ is dominated by $y^3/375$. On
$[\frac{11}5,\frac{13}5]$ the derivative is at most the linear coefficient plus three times the
cubic coefficient, which is negative, so the quartic decreases and its minimum is its value at
$y=1$, namely $\frac{89}{87500}$ for $\phi_A$ and $\frac{9203}{800000}$ for $\phi_B$. The tightest
case is $\phi_A$ at $\nu=13/5$:
\[
\phi_A\bigl(\tfrac{13}5\bigr)=\frac{348293}{1892800}=\frac{23}{125}+\frac{89}{9464000}.
\]
Another argument, which keeps the exact offset $e_0-\varphi$ at the vertex $B$ and minimizes
$\phi_A$ by differentiation, is in \path{HAND_FORWARD_CONDITIONAL_NORM_BOUNDS}, Section~2,
of~\cite{Notes}.

\subsection{The small forward minimum}\label{app:fwd-small}

Suppose $\nu\ge3$ and, after the exchange of sides, $z\le w$. The left row and the floors give
$\mathcal E(h)\ge0$ for the residual
\[
\mathcal E(h)=\eta A-B-n_1F_{z/\hat s_-}(h)-n_2V,\qquad
V=F_{w/\hat s_+}\Bigl(\max\Bigl\{1,\frac{\hat s_+^2}{\hat K_+}\Bigr\}\Bigr),
\]
the floor for $v$ at the least factor that minimality allows. Lowering $w$ does not decrease
$\mathcal E$, by the increments of \cref{app:fwd-centered} and $V\ge\beta(w)$, as long as $q\ge1$.
Lower $w$ until either $q=1$ or $w=z$. Three faces remain:
\begin{enumerate}
\item $q=1$, that is, $w=\hat K/\hat s_--1$ (possible when $\hat s_-^2\le\hat K$);
\item $w=z$ with $\hat s_-^2\ge\hat K$ and $\hat s_-^2\le\hat K_-$;
\item $w=z$ with $\hat s_-^2\ge\hat K$ and $\hat s_-^2>\hat K_-$, the hard face.
\end{enumerate}
On each face $\mathcal E$ has the form in $h$ of \cref{app:fwd-signs}, with second derivative
$-zP(h)/(\hat K_-h^3)$ and the same $P$. Since $P(1)=-z\hat K_-d$ with
$d=\hat s_-^2/\hat K_--1$ and $P'(h)\ge\hat K_-$ for $h\ge1$, integrating twice from $h=1$ with
$\int_1^h(h-r)r^{-3}\,\dd r=(h-1)^2/(2h)$ and $\int_1^h(h-r)(r-1)r^{-3}\,\dd r=f(h)$ gives
\begin{equation}\label{eq:fwd-integral}
\mathcal E(h)\le\mathcal E(1)+(h-1)\,\mathcal E'(1)+\frac{z^2d\,(h-1)^2}{2h}-z\,f(h).
\end{equation}
The value and slope at $h=1$ are explicit: with
$\psi(q)=(q-1)/\hat s_++\frac w{\hat s_+}(1-q^{-1})-\log q$,
\begin{align*}
\mathcal E(1)&=\frac{\hat K}{\hat K_-}\,\psi(q)-\frac{\hat K^2(q-1)}{\hat s_-\hat K_-}\bigl(V-\beta(w)\bigr),\\
\mathcal E'(1)&=\frac{\hat s_-\hat K}{\hat K_-}\Bigl(\psi(q)+\frac{z(q-1)(q-w)}{\hat s_-\hat s_+q}\Bigr)
 -\frac{\hat K^2(\hat s_-q-1)}{\hat s_-\hat K_-}\bigl(V-\beta(w)\bigr).
\end{align*}
Since $\psi'(q)=(q-1)(q-w)/(\hat s_+q^2)$, $\psi$ first decreases and then increases on $[1,\infty)$.
\begin{enumerate}
\item On the face $q=1$: $\psi(1)=0$, and $V=\beta(w)$ because $\hat s_+^2\le\hat K_+$ there; so
  $\mathcal E(1)=\mathcal E'(1)=0$, $d\le0$, and \eqref{eq:fwd-integral} gives
  $\mathcal E(h)\le-zf(h)<0$.
\item On the second face: $V=\beta(z)$, $d\le0$, and $1\le q\le z$, so $\psi(q)\le0$ and the extra
  slope term is nonpositive; again $\mathcal E(h)\le-zf(h)<0$.
\item On the hard face, $\psi(q)\le0$ still holds; by the shape of $\psi$ it suffices to check the
  largest admissible $q$. If $z\le3$, then $\nu\ge3$ gives $q\le\hat s_-^2/4$, and
  $\psi(\hat s_-^2/4)=\frac{\hat s_-}4+1-\frac2{\hat s_-}-\frac4{\hat s_-^2}+\frac4{\hat s_-^3}
  -2\log\frac{\hat s_-}2$ vanishes at $\hat s_-=2$ and has derivative
  $(\hat s_--6)(\hat s_--2)^2(\hat s_-+2)/(4\hat s_-^4)\le0$ for $2\le\hat s_-\le4$. If $z\ge3$, then
  $\nu\ge z$ gives $q\le\hat s_-$, and $\psi(\hat s_-)=x+x^2-\log\hat s_-$ with $x=z/\hat s_-$
  decreases in $\hat s_-$ (its derivative is $-(\hat s_--1)(\hat s_--2)/\hat s_-^3$) from the value
  $\frac{21}{16}-\log4<0$ at $\hat s_-=4$.
  The opposite floor exceeds $\beta(z)$ by at least $\Delta_0=x(x-\frac12)d+xd^2/(3(d+2))$, by the
  Pad\'e bound, and $f(h)\ge(h-1)^3/(3h(h+1))$. Inserting these into \eqref{eq:fwd-integral} gives
  a rational upper bound $\mathcal B$ for $\mathcal E(h)$. Put $\mu=(2z+1-\nu)/\hat s_-\in(0,1)$,
  $\lambda=1-\mu$ and $v=\hat s_-(h-1)/h$, so that $\mu\le v\le1$ and
  $\hat K_-=\hat s_-(\hat s_--\mu)$, $\hat K=\hat s_-(2-\mu)$, $q=\hat s_-/(2-\mu)$,
  $d=\mu/(\hat s_--\mu)$. After the two weakenings
  $\Delta_0\ge z\mu(3(z-1)+\mu)/(6\hat s_-^2(\hat s_--\mu))$ and
  $(h-1)^3/(3h(h+1))\ge(h-1)^3/(6h^2)$, an exact simplification gives
  $\mathcal B\le-zh\,P_5/\bigl(6\hat s_-\hat K_-^2\bigr)$ with the cubic
  \begin{align*}
  P_5={}&(z+\lambda)^2v^3-3z\mu\hat s_-(z+\lambda)v^2\\
   &+\bigl[(1+\lambda)\mu\hat s_-z(3(z-1)+\mu)-6(z+\lambda)(z-\lambda)(1-z\lambda)\bigr]v\\
   &+(1+\lambda)\mu\hat s_-(3(z-1)+\mu)(z-\lambda).
  \end{align*}
  It remains to show $P_5>0$.
\end{enumerate}

\runin{The cubic $P_5$.} We follow \path{HAND_H5_INTEGRAL_RATIONAL_BOUND} of~\cite{Notes}. Put
$\varpi=(2-\mu)z-(2+\mu)=\nu-3\ge0$. Then
$(2-\mu)^3P_5=R_0+R_1\varpi+R_2\varpi^2+R_3\varpi^3$ with
\begin{align*}
R_3={}&3\bigl((2-\mu^2)v+\mu(2-\mu-v^2)\bigr),\\
R_2={}&(2-\mu)v^3-3\mu(\mu^2-\mu+10)v^2\\
 &+(\mu^4-13\mu^3-14\mu^2+24\mu+24)v+E_2,\\
R_1={}&2(2-\mu)(\mu^2-2\mu+4)v^3-3\mu(\mu^3+4\mu^2-4\mu+32)v^2\\
 &+\mu(7\mu^4-46\mu^3+40\mu^2-96\mu+144)v+C_1,\\
R_0={}&(2-\mu)(\mu^2-2\mu+4)^2v^3-12\mu(\mu^3+8)v^2\\
 &+2\mu^2(3\mu^4-16\mu^3+20\mu^2-56\mu+64)v\\
 &+4\mu^3(2-\mu)(4-\mu)(8-\mu),
\end{align*}
where $E_2=4\mu^4+28\mu^2(1-\mu)+24\mu$ and $C_1=\mu^2(-\mu^4+14\mu^3-40\mu^2-48\mu+160)$; here
$0\le\mu\le v\le1$. In $R_1$ and $R_2$ the cubic coefficient is nonnegative and the quadratic one
nonpositive, so $v^2\le v$ gives a bound linear in $v$:
\begin{itemize}
\item $R_3\ge0$, since $2-\mu-v^2\ge1-\mu\ge0$;
\item $R_2\ge(1-v)E_2+v(E_2+D_2)$ with $D_2=\mu^4-16\mu^3-11\mu^2-6\mu+24$, and
  $E_2+D_2=5\mu^4-44\mu^3+17\mu^2+18\mu+24\ge-27\mu^2+18\mu+24=15+9(1-\mu)+27\mu(1-\mu)$;
\item $R_1\ge(1-v)C_1+v(C_1+D_1)$ with $D_1=\mu(7\mu^4-49\mu^3+28\mu^2-84\mu+48)$, where
  $C_1\ge71\mu^2$ and $C_1+D_1=\mu(-\mu^5+21\mu^4-89\mu^3-20\mu^2+76\mu+48)\ge\mu(48-34\mu)
  \ge14\mu$.
\end{itemize}
For $R_0$ there are two cases. If $\mu\le1/5$, its three $v$-coefficients are at least $23$,
at least $-97\mu$ and at least $105\mu^2$, and its constant term is nonnegative, so
$R_0\ge v\bigl(23(v-\frac{97\mu}{46})^2+\frac{251}{92}\mu^2\bigr)>0$ for $v>0$. If $\mu\ge1/5$,
write $v=\mu+(1-\mu)s$ with $s\in[0,1]$; then
\[
R_0=\mu^3A_0(1-s)^3+\mu^2A_1s(1-s)^2+\mu^2A_2s^2(1-s)+Q_0s^3
\]
with
\begin{align*}
A_0&=320-384\mu+136\mu^2-68\mu^3+12\mu^4-\mu^5,\\
A_1&=9\mu^5-64\mu^4+132\mu^3-736\mu^2+672\mu+32,\\
A_2&=6\mu^5-35\mu^4+138\mu^3-776\mu^2+792\mu-80,\\
Q_0&=2\mu^6+23\mu^5-190\mu^4+124\mu^3+168\mu^2-144\mu+32 .
\end{align*}
In $n=1-\mu$, $A_0=n^5+7n^4+30n^3-6n^2+273n+15>0$. The polynomials $A_1,A_2$ are concave on
$[0,1]$ (their second derivatives are $4(45\mu^3-192\mu^2+198\mu-368)<0$ and
$4(30\mu^3-105\mu^2+207\mu-388)<0$), with $A_1(0)=32$, $A_1(1)=45$,
$A_2(\frac15)=\frac{151281}{3125}$ and $A_2(1)=45$. For $\mu\le\frac12$, $-190\mu^4\ge-95\mu^3$ gives
$Q_0\ge29\mu^3+168(\mu-\frac37)^2+\frac87>0$; for $\mu\ge\frac12$, with $n=1-\mu\le\frac12$ and
$r=2n$,
\begin{align*}
Q_0&=2n^6-35n^5-45n^4+366n^3-340n^2+69n+15\\
&\ge15(1-r)^3+\tfrac{159}2r(1-r)^2+29r^2(1-r)+\tfrac{203}{32}r^3>0,
\end{align*}
using $n^5\le n^3/4$ and $n^4\le n^3/2$. So $R_0>0$, hence $P_5>0$, and $\mathcal E(h)<0$ on every
face, a contradiction; therefore $\nu<3$. Another argument for the first two faces, by concavity
in $h$ alone, is in \path{MIXED_GLOBAL_SMALL_FORWARD_MINIMUM}, Sections~4--5, of~\cite{Notes}.
